\documentclass[tikz,border=8pt]{standalone}
\usepackage[UTF8,fontset=windows]{ctex}
\usepackage{tikz}
\usetikzlibrary{arrows.meta, positioning, calc}

\begin{document}
\begin{tikzpicture}[
    font=\small,
    >=Latex,
    line/.style={-Latex, very thick, draw=gray!70},
    dashline/.style={-Latex, thick, draw=gray!55, dashed},
    box/.style={draw, rounded corners=6pt, very thick, align=center, minimum height=11mm},
    srca/.style={box, fill=violet!25, draw=violet!70!black, text width=4.0cm},
    srcb/.style={box, fill=cyan!22, draw=cyan!55!black, text width=4.0cm},
    mux/.style={box, fill=gray!18, draw=gray!65!black, text width=3.6cm},
    proc/.style={box, fill=green!24, draw=green!45!black, text width=5.0cm},
    outbox/.style={box, fill=gray!24, draw=gray!65!black, text width=5.1cm},
    note/.style={draw=none, align=left, font=\scriptsize, text width=4.6cm},
    badge/.style={circle, fill=white, draw=green!45!black, very thick, minimum size=6.5mm, inner sep=0pt, font=\bfseries\small}
]

\node[font=\Large\bfseries] (title) at (0, 1.2) {PDCCH 信道译码流程图};

% source blocks
\node[srca] (src_left) at (-4.8, 0) {上游控制信息\\CFI / 控制区大小\\PCFICH / PHICH 占用信息};
\node[srcb] (src_right) at (4.8, 0) {主输入\\FFT 资源网格\\PCI / SFN / 天线配置};
\node[mux] (mux) at (0, -2.1) {输入汇聚 / DTO 构建\\\texttt{PdcchMmseInput}};

% main chain
\node[proc] (step1) at (0, -4.5) {控制区参数整理与 RE 排除\\构造有效控制区资源\\排除 CRS / PCFICH / PHICH};
\node[badge] (n1) at (-2.8, -4.5) {1};

\node[proc] (step2) at (0, -6.9) {信道估计与 MMSE 均衡\\CRS 信道估计 $\rightarrow$ PDCCH 有效 RE 软符号};
\node[badge] (n2) at (-2.8, -6.9) {2};
\node[note] (note2) at (4.9, -6.9) {工程实现输出：\\$x\_hat\_{re}$、$x\_hat\_{im}$、$sinr$\\以及 RE 来源索引};

\node[proc] (step3) at (0, -9.4) {发射分集去映射（2Tx 可选）\\1Tx：逐 RE 直出\\2Tx：按 RE 对恢复两个 QPSK 软符号};
\node[badge] (n3) at (-2.8, -9.4) {3};
\node[note] (note3) at (4.9, -9.4) {本工程新增：\\\texttt{run\_pdcch\_td(...)}\\输出 \texttt{re\_grid\_indices0/1}};

\node[proc] (step4) at (0, -11.9) {QPSK 软比特 / LLR\\软符号 $\rightarrow$ 每比特置信度};
\node[badge] (n4) at (-2.8, -11.9) {4};

\node[proc] (step5) at (0, -14.4) {REG / CCE 重组 + 搜索空间候选\\聚合级别、候选位置、盲检索路径};
\node[badge] (n5) at (-2.8, -14.4) {5};

\node[proc] (step6) at (0, -16.9) {候选译码\\解扰 $\rightarrow$ 速率恢复\\卷积译码 $\rightarrow$ CRC-RNTI};
\node[badge] (n6) at (-2.8, -16.9) {6};
\node[note] (note6) at (4.9, -16.9) {标准完整接收机需要：\\盲检索全部候选并做 RNTI 掩码 CRC 校验};

\node[outbox] (step7) at (0, -19.6) {DCI 输出\\RNTI / DCI 格式 / CCE / 聚合级别\\调度授权 / HARQ / MCS / RV / NDI};
\node[badge] (n7) at (-2.8, -19.6) {7};

% arrows
\draw[line] (src_left) -- ($(mux.north west)+(0.3,0)$);
\draw[line] (src_right) -- ($(mux.north east)+(-0.3,0)$);
\draw[line] (mux) -- (step1);
\draw[line] (step1) -- (step2);
\draw[line] (step2) -- (step3);
\draw[line] (step3) -- (step4);
\draw[line] (step4) -- (step5);
\draw[line] (step5) -- (step6);
\draw[line] (step6) -- (step7);

% footer note
\node[note, text width=11.0cm] at (0, -22.3) {说明：当前 \texttt{MMSE\_CPP} 主要覆盖步骤 1--3，即控制区资源整理、CRS 信道估计、MMSE 均衡与 2Tx 发射分集去映射；步骤 4--7 属于完整 PDCCH 接收机后续链路，通常在仓库外部完成。};

\end{tikzpicture}
\end{document}
